% Visual layer on top of common.tex: tiles, icons, a progress-strip section divider, and a footer
% strip that defines the concepts appearing for the first time on each slide.
% Same look as the Class 20 deck. Load with:  \documentclass[aspectratio=169]{beamer}

\usetheme{Boadilla}
\usefonttheme{professionalfonts}
\setbeamertemplate{navigation symbols}{}
\setbeamertemplate{itemize items}[circle]
\setbeamertemplate{enumerate items}[default]
\setbeamertemplate{headline}{}

\usepackage[utf8]{inputenc}
\usepackage[T1]{fontenc}
\usepackage{lmodern}
\usepackage{amsmath}
\usepackage{booktabs}
\usepackage{tabularx}
\usepackage{array}
\usepackage{hyperref}
\usepackage[strings]{underscore}
\usepackage{listings}
\usepackage{tikz}
\usetikzlibrary{arrows.meta,positioning,fit,calc,shapes.geometric}
\usepackage{xcolor}

\definecolor{courseblue}{HTML}{12355B}
\definecolor{coursegold}{HTML}{C98A2E}
\definecolor{courseice}{HTML}{EEF3F8}
\definecolor{coursegreen}{HTML}{2F6B4F}
\definecolor{coursegray}{HTML}{4B5563}
\definecolor{codebg}{HTML}{F7F9FC}
\definecolor{codecomment}{HTML}{5E6C84}
\definecolor{codestring}{HTML}{0B6E4F}

\setbeamercolor{palette primary}{bg=courseblue,fg=white}
\setbeamercolor{palette secondary}{bg=coursegold,fg=white}
\setbeamercolor{palette tertiary}{bg=courseblue,fg=white}
\setbeamercolor{palette quaternary}{bg=coursegold,fg=white}
\setbeamercolor{structure}{fg=courseblue}
\setbeamercolor{frametitle}{bg=courseice,fg=courseblue}
\setbeamercolor{title}{fg=courseblue}
\setbeamercolor{block title}{bg=courseblue,fg=white}
\setbeamercolor{block body}{bg=courseice,fg=black}
\setbeamercolor{normal text}{fg=black,bg=white}
\setbeamercolor{alerted text}{fg=coursegold}

\setbeamerfont{footline}{size=\scriptsize}

\setbeamertemplate{footline}{
  \leavevmode
  \hbox{
    \begin{beamercolorbox}[wd=.33\paperwidth,ht=2.8ex,dp=1.2ex,leftskip=0.8em]{palette primary}%
      University of Trieste
    \end{beamercolorbox}%
    \begin{beamercolorbox}[wd=.49\paperwidth,ht=2.8ex,dp=1.2ex,center]{palette tertiary}%
      \coursename{} \textbar{} \coursecode
    \end{beamercolorbox}%
    \begin{beamercolorbox}[wd=.18\paperwidth,ht=2.8ex,dp=1.2ex,rightskip=0.8em plus1fil]{palette secondary}%
      \insertframenumber{} / \inserttotalframenumber
    \end{beamercolorbox}%
  }
}

\hypersetup{
  colorlinks=true,
  linkcolor=courseblue,
  urlcolor=coursegreen
}

\lstdefinestyle{coursepython}{
  language=Python,
  basicstyle=\ttfamily\footnotesize,
  keywordstyle=\color{courseblue}\bfseries,
  commentstyle=\color{codecomment}\itshape,
  stringstyle=\color{codestring},
  showstringspaces=false,
  columns=fullflexible,
  keepspaces=true,
  frame=single,
  framerule=0.6pt,
  rulecolor=\color{courseblue},
  backgroundcolor=\color{codebg},
  xleftmargin=0.4em,
  xrightmargin=0.4em,
  aboveskip=0.4em,
  belowskip=0.2em
}

\newcommand{\coursename}{Data-driven Systems Engineering}
\newcommand{\coursecode}{440MI / 305SM}
\newcommand{\coursefooter}{}

\newcommand{\learninggoals}[1]{
  \begin{block}{Learning goals}
  #1
  \end{block}
}

\newcommand{\repoactivity}[1]{
  \begin{block}{Repo activity}
  #1
  \end{block}
}

\newcommand{\readinglist}[1]{
  \begin{block}{Papers and materials}
  {\footnotesize #1}
  \end{block}
}

\newcommand{\coursecontextframe}[5]{
  \begin{frame}{Course Map and Running Cases}
  \centering
  \begin{tikzpicture}[node distance=0.28cm]
  \node[flowbox, minimum width=1.45cm, minimum height=0.8cm] (problem) {\scriptsize Problem};
  \node[flowbox, minimum width=1.45cm, minimum height=0.8cm, right=of problem] (data) {\scriptsize Data};
  \node[flowbox, minimum width=1.45cm, minimum height=0.8cm, right=of data] (features) {\scriptsize Features};
  \node[flowbox, minimum width=1.45cm, minimum height=0.8cm, right=of features] (model) {\scriptsize Model};
  \node[flowbox, minimum width=1.45cm, minimum height=0.8cm, right=of model] (interface) {\scriptsize Interface};
  \node[flowbox, minimum width=1.45cm, minimum height=0.8cm, right=of interface] (delivery) {\scriptsize Delivery};
  \node[flowbox, minimum width=1.45cm, minimum height=0.8cm, right=of delivery] (monitor) {\scriptsize Monitor};
  \node[accentbox, minimum width=1.6cm, minimum height=0.8cm, right=of monitor] (evolution) {\scriptsize Evolution};
  \foreach \a/\b in {problem/data,data/features,features/model,model/interface,interface/delivery,delivery/monitor,monitor/evolution} {
    \draw[arrow] (\a) -- (\b);
  }
  \end{tikzpicture}

  \vspace{0.5em}
  \begin{block}{Today in the course}
  \textbf{Focus:} #1

  \textbf{Bridge:} #2
  \end{block}

  \begin{columns}[T]
  \column{0.48\textwidth}
  \begin{block}{Fraud case}
  {\footnotesize #3}
  \end{block}
  \column{0.48\textwidth}
  \begin{block}{Pasteurization case}
  {\footnotesize #4}
  \end{block}
  \end{columns}

  \vspace{0.3em}
  {\footnotesize \textbf{Course thread:} #5}
  \coursefooter
  \end{frame}
}

\newcommand{\classclosingframe}[4]{
  \begin{frame}{From This Class to the Next}
  \begin{block}{What stays with us}
  #1
  \end{block}

  \begin{columns}[T]
  \column{0.48\textwidth}
  \begin{block}{Project use}
  {\footnotesize #2}
  \end{block}
  \column{0.48\textwidth}
  \begin{block}{Next connection}
  {\footnotesize #3}
  \end{block}
  \end{columns}

  \vspace{0.3em}
  {\footnotesize \textbf{Course thread:} #4}
  \coursefooter
  \end{frame}
}

\tikzset{
  flowbox/.style={
    rounded corners=3pt,
    draw=courseblue,
    fill=courseice,
    very thick,
    minimum width=2.2cm,
    minimum height=0.9cm,
    align=center,
    font=\small
  },
  accentbox/.style={
    rounded corners=3pt,
    draw=coursegold,
    fill=coursegold!15,
    very thick,
    minimum width=2.2cm,
    minimum height=0.9cm,
    align=center,
    font=\small
  },
  arrow/.style={
    -{Latex[length=2.5mm,width=2mm]},
    thick,
    draw=courseblue
  }
}
  % Visual layer on top of common.tex: tiles, icons, a progress-strip section divider, and a footer
% strip that defines the concepts appearing for the first time on each slide.
% Same look as the Class 20 deck. Load with:  \documentclass[aspectratio=169]{beamer}

\usetheme{Boadilla}
\usefonttheme{professionalfonts}
\setbeamertemplate{navigation symbols}{}
\setbeamertemplate{itemize items}[circle]
\setbeamertemplate{enumerate items}[default]
\setbeamertemplate{headline}{}

\usepackage[utf8]{inputenc}
\usepackage[T1]{fontenc}
\usepackage{lmodern}
\usepackage{amsmath}
\usepackage{booktabs}
\usepackage{tabularx}
\usepackage{array}
\usepackage{hyperref}
\usepackage[strings]{underscore}
\usepackage{listings}
\usepackage{tikz}
\usetikzlibrary{arrows.meta,positioning,fit,calc,shapes.geometric}
\usepackage{xcolor}

\definecolor{courseblue}{HTML}{12355B}
\definecolor{coursegold}{HTML}{C98A2E}
\definecolor{courseice}{HTML}{EEF3F8}
\definecolor{coursegreen}{HTML}{2F6B4F}
\definecolor{coursegray}{HTML}{4B5563}
\definecolor{codebg}{HTML}{F7F9FC}
\definecolor{codecomment}{HTML}{5E6C84}
\definecolor{codestring}{HTML}{0B6E4F}

\setbeamercolor{palette primary}{bg=courseblue,fg=white}
\setbeamercolor{palette secondary}{bg=coursegold,fg=white}
\setbeamercolor{palette tertiary}{bg=courseblue,fg=white}
\setbeamercolor{palette quaternary}{bg=coursegold,fg=white}
\setbeamercolor{structure}{fg=courseblue}
\setbeamercolor{frametitle}{bg=courseice,fg=courseblue}
\setbeamercolor{title}{fg=courseblue}
\setbeamercolor{block title}{bg=courseblue,fg=white}
\setbeamercolor{block body}{bg=courseice,fg=black}
\setbeamercolor{normal text}{fg=black,bg=white}
\setbeamercolor{alerted text}{fg=coursegold}

\setbeamerfont{footline}{size=\scriptsize}

\setbeamertemplate{footline}{
  \leavevmode
  \hbox{
    \begin{beamercolorbox}[wd=.33\paperwidth,ht=2.8ex,dp=1.2ex,leftskip=0.8em]{palette primary}%
      University of Trieste
    \end{beamercolorbox}%
    \begin{beamercolorbox}[wd=.49\paperwidth,ht=2.8ex,dp=1.2ex,center]{palette tertiary}%
      \coursename{} \textbar{} \coursecode
    \end{beamercolorbox}%
    \begin{beamercolorbox}[wd=.18\paperwidth,ht=2.8ex,dp=1.2ex,rightskip=0.8em plus1fil]{palette secondary}%
      \insertframenumber{} / \inserttotalframenumber
    \end{beamercolorbox}%
  }
}

\hypersetup{
  colorlinks=true,
  linkcolor=courseblue,
  urlcolor=coursegreen
}

\lstdefinestyle{coursepython}{
  language=Python,
  basicstyle=\ttfamily\footnotesize,
  keywordstyle=\color{courseblue}\bfseries,
  commentstyle=\color{codecomment}\itshape,
  stringstyle=\color{codestring},
  showstringspaces=false,
  columns=fullflexible,
  keepspaces=true,
  frame=single,
  framerule=0.6pt,
  rulecolor=\color{courseblue},
  backgroundcolor=\color{codebg},
  xleftmargin=0.4em,
  xrightmargin=0.4em,
  aboveskip=0.4em,
  belowskip=0.2em
}

\newcommand{\coursename}{Data-driven Systems Engineering}
\newcommand{\coursecode}{440MI / 305SM}
\newcommand{\coursefooter}{}

\newcommand{\learninggoals}[1]{
  \begin{block}{Learning goals}
  #1
  \end{block}
}

\newcommand{\repoactivity}[1]{
  \begin{block}{Repo activity}
  #1
  \end{block}
}

\newcommand{\readinglist}[1]{
  \begin{block}{Papers and materials}
  {\footnotesize #1}
  \end{block}
}

\newcommand{\coursecontextframe}[5]{
  \begin{frame}{Course Map and Running Cases}
  \centering
  \begin{tikzpicture}[node distance=0.28cm]
  \node[flowbox, minimum width=1.45cm, minimum height=0.8cm] (problem) {\scriptsize Problem};
  \node[flowbox, minimum width=1.45cm, minimum height=0.8cm, right=of problem] (data) {\scriptsize Data};
  \node[flowbox, minimum width=1.45cm, minimum height=0.8cm, right=of data] (features) {\scriptsize Features};
  \node[flowbox, minimum width=1.45cm, minimum height=0.8cm, right=of features] (model) {\scriptsize Model};
  \node[flowbox, minimum width=1.45cm, minimum height=0.8cm, right=of model] (interface) {\scriptsize Interface};
  \node[flowbox, minimum width=1.45cm, minimum height=0.8cm, right=of interface] (delivery) {\scriptsize Delivery};
  \node[flowbox, minimum width=1.45cm, minimum height=0.8cm, right=of delivery] (monitor) {\scriptsize Monitor};
  \node[accentbox, minimum width=1.6cm, minimum height=0.8cm, right=of monitor] (evolution) {\scriptsize Evolution};
  \foreach \a/\b in {problem/data,data/features,features/model,model/interface,interface/delivery,delivery/monitor,monitor/evolution} {
    \draw[arrow] (\a) -- (\b);
  }
  \end{tikzpicture}

  \vspace{0.5em}
  \begin{block}{Today in the course}
  \textbf{Focus:} #1

  \textbf{Bridge:} #2
  \end{block}

  \begin{columns}[T]
  \column{0.48\textwidth}
  \begin{block}{Fraud case}
  {\footnotesize #3}
  \end{block}
  \column{0.48\textwidth}
  \begin{block}{Pasteurization case}
  {\footnotesize #4}
  \end{block}
  \end{columns}

  \vspace{0.3em}
  {\footnotesize \textbf{Course thread:} #5}
  \coursefooter
  \end{frame}
}

\newcommand{\classclosingframe}[4]{
  \begin{frame}{From This Class to the Next}
  \begin{block}{What stays with us}
  #1
  \end{block}

  \begin{columns}[T]
  \column{0.48\textwidth}
  \begin{block}{Project use}
  {\footnotesize #2}
  \end{block}
  \column{0.48\textwidth}
  \begin{block}{Next connection}
  {\footnotesize #3}
  \end{block}
  \end{columns}

  \vspace{0.3em}
  {\footnotesize \textbf{Course thread:} #4}
  \coursefooter
  \end{frame}
}

\tikzset{
  flowbox/.style={
    rounded corners=3pt,
    draw=courseblue,
    fill=courseice,
    very thick,
    minimum width=2.2cm,
    minimum height=0.9cm,
    align=center,
    font=\small
  },
  accentbox/.style={
    rounded corners=3pt,
    draw=coursegold,
    fill=coursegold!15,
    very thick,
    minimum width=2.2cm,
    minimum height=0.9cm,
    align=center,
    font=\small
  },
  arrow/.style={
    -{Latex[length=2.5mm,width=2mm]},
    thick,
    draw=courseblue
  }
}
  % Visual layer on top of common.tex: tiles, icons, a progress-strip section divider, and a footer
% strip that defines the concepts appearing for the first time on each slide.
% Same look as the Class 20 deck. Load with:  \documentclass[aspectratio=169]{beamer}

\usetheme{Boadilla}
\usefonttheme{professionalfonts}
\setbeamertemplate{navigation symbols}{}
\setbeamertemplate{itemize items}[circle]
\setbeamertemplate{enumerate items}[default]
\setbeamertemplate{headline}{}

\usepackage[utf8]{inputenc}
\usepackage[T1]{fontenc}
\usepackage{lmodern}
\usepackage{amsmath}
\usepackage{booktabs}
\usepackage{tabularx}
\usepackage{array}
\usepackage{hyperref}
\usepackage[strings]{underscore}
\usepackage{listings}
\usepackage{tikz}
\usetikzlibrary{arrows.meta,positioning,fit,calc,shapes.geometric}
\usepackage{xcolor}

\definecolor{courseblue}{HTML}{12355B}
\definecolor{coursegold}{HTML}{C98A2E}
\definecolor{courseice}{HTML}{EEF3F8}
\definecolor{coursegreen}{HTML}{2F6B4F}
\definecolor{coursegray}{HTML}{4B5563}
\definecolor{codebg}{HTML}{F7F9FC}
\definecolor{codecomment}{HTML}{5E6C84}
\definecolor{codestring}{HTML}{0B6E4F}

\setbeamercolor{palette primary}{bg=courseblue,fg=white}
\setbeamercolor{palette secondary}{bg=coursegold,fg=white}
\setbeamercolor{palette tertiary}{bg=courseblue,fg=white}
\setbeamercolor{palette quaternary}{bg=coursegold,fg=white}
\setbeamercolor{structure}{fg=courseblue}
\setbeamercolor{frametitle}{bg=courseice,fg=courseblue}
\setbeamercolor{title}{fg=courseblue}
\setbeamercolor{block title}{bg=courseblue,fg=white}
\setbeamercolor{block body}{bg=courseice,fg=black}
\setbeamercolor{normal text}{fg=black,bg=white}
\setbeamercolor{alerted text}{fg=coursegold}

\setbeamerfont{footline}{size=\scriptsize}

\setbeamertemplate{footline}{
  \leavevmode
  \hbox{
    \begin{beamercolorbox}[wd=.33\paperwidth,ht=2.8ex,dp=1.2ex,leftskip=0.8em]{palette primary}%
      University of Trieste
    \end{beamercolorbox}%
    \begin{beamercolorbox}[wd=.49\paperwidth,ht=2.8ex,dp=1.2ex,center]{palette tertiary}%
      \coursename{} \textbar{} \coursecode
    \end{beamercolorbox}%
    \begin{beamercolorbox}[wd=.18\paperwidth,ht=2.8ex,dp=1.2ex,rightskip=0.8em plus1fil]{palette secondary}%
      \insertframenumber{} / \inserttotalframenumber
    \end{beamercolorbox}%
  }
}

\hypersetup{
  colorlinks=true,
  linkcolor=courseblue,
  urlcolor=coursegreen
}

\lstdefinestyle{coursepython}{
  language=Python,
  basicstyle=\ttfamily\footnotesize,
  keywordstyle=\color{courseblue}\bfseries,
  commentstyle=\color{codecomment}\itshape,
  stringstyle=\color{codestring},
  showstringspaces=false,
  columns=fullflexible,
  keepspaces=true,
  frame=single,
  framerule=0.6pt,
  rulecolor=\color{courseblue},
  backgroundcolor=\color{codebg},
  xleftmargin=0.4em,
  xrightmargin=0.4em,
  aboveskip=0.4em,
  belowskip=0.2em
}

\newcommand{\coursename}{Data-driven Systems Engineering}
\newcommand{\coursecode}{440MI / 305SM}
\newcommand{\coursefooter}{}

\newcommand{\learninggoals}[1]{
  \begin{block}{Learning goals}
  #1
  \end{block}
}

\newcommand{\repoactivity}[1]{
  \begin{block}{Repo activity}
  #1
  \end{block}
}

\newcommand{\readinglist}[1]{
  \begin{block}{Papers and materials}
  {\footnotesize #1}
  \end{block}
}

\newcommand{\coursecontextframe}[5]{
  \begin{frame}{Course Map and Running Cases}
  \centering
  \begin{tikzpicture}[node distance=0.28cm]
  \node[flowbox, minimum width=1.45cm, minimum height=0.8cm] (problem) {\scriptsize Problem};
  \node[flowbox, minimum width=1.45cm, minimum height=0.8cm, right=of problem] (data) {\scriptsize Data};
  \node[flowbox, minimum width=1.45cm, minimum height=0.8cm, right=of data] (features) {\scriptsize Features};
  \node[flowbox, minimum width=1.45cm, minimum height=0.8cm, right=of features] (model) {\scriptsize Model};
  \node[flowbox, minimum width=1.45cm, minimum height=0.8cm, right=of model] (interface) {\scriptsize Interface};
  \node[flowbox, minimum width=1.45cm, minimum height=0.8cm, right=of interface] (delivery) {\scriptsize Delivery};
  \node[flowbox, minimum width=1.45cm, minimum height=0.8cm, right=of delivery] (monitor) {\scriptsize Monitor};
  \node[accentbox, minimum width=1.6cm, minimum height=0.8cm, right=of monitor] (evolution) {\scriptsize Evolution};
  \foreach \a/\b in {problem/data,data/features,features/model,model/interface,interface/delivery,delivery/monitor,monitor/evolution} {
    \draw[arrow] (\a) -- (\b);
  }
  \end{tikzpicture}

  \vspace{0.5em}
  \begin{block}{Today in the course}
  \textbf{Focus:} #1

  \textbf{Bridge:} #2
  \end{block}

  \begin{columns}[T]
  \column{0.48\textwidth}
  \begin{block}{Fraud case}
  {\footnotesize #3}
  \end{block}
  \column{0.48\textwidth}
  \begin{block}{Pasteurization case}
  {\footnotesize #4}
  \end{block}
  \end{columns}

  \vspace{0.3em}
  {\footnotesize \textbf{Course thread:} #5}
  \coursefooter
  \end{frame}
}

\newcommand{\classclosingframe}[4]{
  \begin{frame}{From This Class to the Next}
  \begin{block}{What stays with us}
  #1
  \end{block}

  \begin{columns}[T]
  \column{0.48\textwidth}
  \begin{block}{Project use}
  {\footnotesize #2}
  \end{block}
  \column{0.48\textwidth}
  \begin{block}{Next connection}
  {\footnotesize #3}
  \end{block}
  \end{columns}

  \vspace{0.3em}
  {\footnotesize \textbf{Course thread:} #4}
  \coursefooter
  \end{frame}
}

\tikzset{
  flowbox/.style={
    rounded corners=3pt,
    draw=courseblue,
    fill=courseice,
    very thick,
    minimum width=2.2cm,
    minimum height=0.9cm,
    align=center,
    font=\small
  },
  accentbox/.style={
    rounded corners=3pt,
    draw=coursegold,
    fill=coursegold!15,
    very thick,
    minimum width=2.2cm,
    minimum height=0.9cm,
    align=center,
    font=\small
  },
  arrow/.style={
    -{Latex[length=2.5mm,width=2mm]},
    thick,
    draw=courseblue
  }
}
  % Visual layer on top of common.tex: tiles, icons, a progress-strip section divider, and a footer
% strip that defines the concepts appearing for the first time on each slide.
% Same look as the Class 20 deck. Load with:  \input{common}  \input{common_visual}
%
% Before the first \sectiondivider, a deck must define
%   \sectioncount   the number of sections, e.g. 4
%   \sectionlabels  a \foreach list of number/label pairs, e.g. 1/Need,2/Testable,3/Manage,4/Practice
% Icons come from the fontawesome5 package (Font Awesome 5 Free), e.g. \faIcon{robot}.

\usepackage{colortbl}
\usepackage{fontawesome5}
\usetikzlibrary{backgrounds}

% Slides read better without automatic hyphenation.
\hyphenpenalty=10000
\exhyphenpenalty=10000

\definecolor{coursered}{HTML}{A23B2A}
\newcommand{\goldink}{\color{coursegold!65!black}}

\tikzset{
  tile/.style={rounded corners=4pt, fill=courseice, draw=none, align=center, font=\footnotesize, inner sep=5pt},
  tileblue/.style={tile, fill=white, draw=courseblue!55, thick},
  tilegold/.style={tile, fill=coursegold!14, draw=coursegold, thick},
  tilegreen/.style={tile, fill=coursegreen!10, draw=coursegreen!70, thick},
  tilered/.style={tile, fill=coursered!8, draw=coursered!70, thick},
  tilegray/.style={tile, fill=white, draw=coursegray!55, thick, dashed},
  tag/.style={rounded corners=2pt, fill=coursegold, text=white, font=\bfseries\tiny, inner sep=2.5pt},
  chip/.style={rounded corners=3pt, fill=courseblue, text=white, font=\footnotesize\bfseries, align=center, inner sep=4pt},
  note/.style={font=\scriptsize, text=coursegray, align=left},
  arrowgold/.style={-{Latex[length=2.5mm,width=2mm]}, thick, draw=coursegold},
  arrowred/.style={-{Latex[length=2.5mm,width=2mm]}, thick, draw=coursered},
  arrowgreen/.style={-{Latex[length=2.5mm,width=2mm]}, thick, draw=coursegreen},
  arrowgray/.style={-{Latex[length=2mm,width=1.6mm]}, thick, draw=coursegray!70},
}

% A round icon badge, drawn inside a tikzpicture: x y colour icon
\newcommand{\badge}[4]{%
  \node[circle, fill=#3, text=white, minimum size=0.85cm, inner sep=0pt] at (#1,#2) {\large\faIcon{#4}};
}

% Inline tags for the kind of requirement.
\newcommand{\Ftag}{\colorbox{courseblue}{\color{white}\bfseries\scriptsize\ \ F\ \ }}
\newcommand{\NFtag}{\colorbox{coursegold}{\color{white}\bfseries\scriptsize\ NF\ }}
\newcommand{\Dtag}{\colorbox{coursegreen}{\color{white}\bfseries\scriptsize\ \ D\ \ }}
\newcommand{\MLtag}{\colorbox{coursered}{\color{white}\bfseries\scriptsize\ ML\ }}

% Footer: a strip with the concepts that appear for the first time on the slide, above the
% usual bar. The bar total counts only the main lecture; a backup slide is numbered separately.
\makeatletter
\def\conceptstext{}
\def\currentsectionname{}
\setbeamercolor{conceptbar}{bg=coursegold!12,fg=black}
\newcommand{\cterm}[2]{{\bfseries\color{courseblue}#1:}\ #2\hspace{0.8em}}
\newcommand{\newtag}{\colorbox{coursegold}{\color{white}\bfseries\tiny\ NEW\ }\hspace{0.6em}}
\setbeamertemplate{footline}{
  \leavevmode
  \vbox{%
    \ifx\conceptstext\empty\else
      \hbox{\color{coursegold}\rule{\paperwidth}{0.8pt}}%
      \hbox{\begin{beamercolorbox}[wd=\paperwidth,leftskip=0.5cm,rightskip=0.5cm,sep=3pt]{conceptbar}\scriptsize\newtag\conceptstext\end{beamercolorbox}}%
    \fi
    \hbox{%
      \begin{beamercolorbox}[wd=.33\paperwidth,ht=2.8ex,dp=1.2ex,leftskip=0.8em]{palette primary}%
        \ifbeamer@inappendix Backup\else\ifnum\value{section}>0 \arabic{section}\ \textbar\ \currentsectionname\else University of Trieste\fi\fi
      \end{beamercolorbox}%
      \begin{beamercolorbox}[wd=.49\paperwidth,ht=2.8ex,dp=1.2ex,center]{palette tertiary}%
        \coursename{} \textbar{} \coursecode
      \end{beamercolorbox}%
      \begin{beamercolorbox}[wd=.18\paperwidth,ht=2.8ex,dp=1.2ex,right,rightskip=0.8em]{palette secondary}%
        \ifbeamer@inappendix Backup \insertframenumberinappendix\else\insertframenumber{} / \insertmainframenumber\fi
      \end{beamercolorbox}%
    }}}
\makeatother
\newenvironment{cframe}[2]{\gdef\conceptstext{#2}\begin{frame}{#1}}{\end{frame}\gdef\conceptstext{}}

% Section divider: progress strip, recap of the previous part, and the roadmap of the next one.
% Arguments: number, title, icon, previously, items.
\newcommand{\sectiondivider}[5]{%
  \section{#2}\gdef\currentsectionname{#2}
  \begin{frame}
  \centering
  \begin{tikzpicture}
  \pgfmathsetmacro{\chipw}{13.6/\sectioncount}
  \foreach \i/\lab in \sectionlabels {
    \ifnum\i<#1 \def\chipfill{courseblue!22}\def\chiptext{courseblue}\else
    \ifnum\i=#1 \def\chipfill{courseblue}\def\chiptext{white}\else
    \def\chipfill{courseice}\def\chiptext{coursegray}\fi\fi
    \node[rounded corners=3pt, minimum width={(\chipw-0.1)*1cm}, minimum height=0.6cm, font=\scriptsize\bfseries, fill=\chipfill, text=\chiptext] at ({\chipw/2+(\i-1)*\chipw},0) {\i\ \lab};
  }
  \node[circle, fill=courseblue, text=white, font=\Huge\bfseries, minimum size=1.7cm, inner sep=0] at (1.0,-1.7) {#1};
  \node[anchor=west, font=\huge\bfseries, text=courseblue, text width=8.8cm, align=left] at (2.2,-1.7) {#2};
  \node[circle, fill=coursegold!16, draw=coursegold, thick, text=coursegold!75!black, font=\fontsize{26}{26}\selectfont, minimum size=1.9cm, inner sep=0] at (12.6,-1.7) {\faIcon{#3}};
  \ifx\relax#4\relax\else
    \node[anchor=north west, font=\small\itshape, text=coursegray, text width=13.2cm, align=left] at (0.1,-2.95) {Previously: #4};
  \fi
  \node[anchor=north west, font=\footnotesize\bfseries, text=coursegold!65!black] at (0.1,-3.95) {IN THIS PART};
  \node[anchor=north west, font=\small, text width=13.2cm, align=left] at (0.1,-4.35) {#5};
  \end{tikzpicture}
  \end{frame}
}

%
% Before the first \sectiondivider, a deck must define
%   \sectioncount   the number of sections, e.g. 4
%   \sectionlabels  a \foreach list of number/label pairs, e.g. 1/Need,2/Testable,3/Manage,4/Practice
% Icons come from the fontawesome5 package (Font Awesome 5 Free), e.g. \faIcon{robot}.

\usepackage{colortbl}
\usepackage{fontawesome5}
\usetikzlibrary{backgrounds}

% Slides read better without automatic hyphenation.
\hyphenpenalty=10000
\exhyphenpenalty=10000

\definecolor{coursered}{HTML}{A23B2A}
\newcommand{\goldink}{\color{coursegold!65!black}}

\tikzset{
  tile/.style={rounded corners=4pt, fill=courseice, draw=none, align=center, font=\footnotesize, inner sep=5pt},
  tileblue/.style={tile, fill=white, draw=courseblue!55, thick},
  tilegold/.style={tile, fill=coursegold!14, draw=coursegold, thick},
  tilegreen/.style={tile, fill=coursegreen!10, draw=coursegreen!70, thick},
  tilered/.style={tile, fill=coursered!8, draw=coursered!70, thick},
  tilegray/.style={tile, fill=white, draw=coursegray!55, thick, dashed},
  tag/.style={rounded corners=2pt, fill=coursegold, text=white, font=\bfseries\tiny, inner sep=2.5pt},
  chip/.style={rounded corners=3pt, fill=courseblue, text=white, font=\footnotesize\bfseries, align=center, inner sep=4pt},
  note/.style={font=\scriptsize, text=coursegray, align=left},
  arrowgold/.style={-{Latex[length=2.5mm,width=2mm]}, thick, draw=coursegold},
  arrowred/.style={-{Latex[length=2.5mm,width=2mm]}, thick, draw=coursered},
  arrowgreen/.style={-{Latex[length=2.5mm,width=2mm]}, thick, draw=coursegreen},
  arrowgray/.style={-{Latex[length=2mm,width=1.6mm]}, thick, draw=coursegray!70},
}

% A round icon badge, drawn inside a tikzpicture: x y colour icon
\newcommand{\badge}[4]{%
  \node[circle, fill=#3, text=white, minimum size=0.85cm, inner sep=0pt] at (#1,#2) {\large\faIcon{#4}};
}

% Inline tags for the kind of requirement.
\newcommand{\Ftag}{\colorbox{courseblue}{\color{white}\bfseries\scriptsize\ \ F\ \ }}
\newcommand{\NFtag}{\colorbox{coursegold}{\color{white}\bfseries\scriptsize\ NF\ }}
\newcommand{\Dtag}{\colorbox{coursegreen}{\color{white}\bfseries\scriptsize\ \ D\ \ }}
\newcommand{\MLtag}{\colorbox{coursered}{\color{white}\bfseries\scriptsize\ ML\ }}

% Footer: a strip with the concepts that appear for the first time on the slide, above the
% usual bar. The bar total counts only the main lecture; a backup slide is numbered separately.
\makeatletter
\def\conceptstext{}
\def\currentsectionname{}
\setbeamercolor{conceptbar}{bg=coursegold!12,fg=black}
\newcommand{\cterm}[2]{{\bfseries\color{courseblue}#1:}\ #2\hspace{0.8em}}
\newcommand{\newtag}{\colorbox{coursegold}{\color{white}\bfseries\tiny\ NEW\ }\hspace{0.6em}}
\setbeamertemplate{footline}{
  \leavevmode
  \vbox{%
    \ifx\conceptstext\empty\else
      \hbox{\color{coursegold}\rule{\paperwidth}{0.8pt}}%
      \hbox{\begin{beamercolorbox}[wd=\paperwidth,leftskip=0.5cm,rightskip=0.5cm,sep=3pt]{conceptbar}\scriptsize\newtag\conceptstext\end{beamercolorbox}}%
    \fi
    \hbox{%
      \begin{beamercolorbox}[wd=.33\paperwidth,ht=2.8ex,dp=1.2ex,leftskip=0.8em]{palette primary}%
        \ifbeamer@inappendix Backup\else\ifnum\value{section}>0 \arabic{section}\ \textbar\ \currentsectionname\else University of Trieste\fi\fi
      \end{beamercolorbox}%
      \begin{beamercolorbox}[wd=.49\paperwidth,ht=2.8ex,dp=1.2ex,center]{palette tertiary}%
        \coursename{} \textbar{} \coursecode
      \end{beamercolorbox}%
      \begin{beamercolorbox}[wd=.18\paperwidth,ht=2.8ex,dp=1.2ex,right,rightskip=0.8em]{palette secondary}%
        \ifbeamer@inappendix Backup \insertframenumberinappendix\else\insertframenumber{} / \insertmainframenumber\fi
      \end{beamercolorbox}%
    }}}
\makeatother
\newenvironment{cframe}[2]{\gdef\conceptstext{#2}\begin{frame}{#1}}{\end{frame}\gdef\conceptstext{}}

% Section divider: progress strip, recap of the previous part, and the roadmap of the next one.
% Arguments: number, title, icon, previously, items.
\newcommand{\sectiondivider}[5]{%
  \section{#2}\gdef\currentsectionname{#2}
  \begin{frame}
  \centering
  \begin{tikzpicture}
  \pgfmathsetmacro{\chipw}{13.6/\sectioncount}
  \foreach \i/\lab in \sectionlabels {
    \ifnum\i<#1 \def\chipfill{courseblue!22}\def\chiptext{courseblue}\else
    \ifnum\i=#1 \def\chipfill{courseblue}\def\chiptext{white}\else
    \def\chipfill{courseice}\def\chiptext{coursegray}\fi\fi
    \node[rounded corners=3pt, minimum width={(\chipw-0.1)*1cm}, minimum height=0.6cm, font=\scriptsize\bfseries, fill=\chipfill, text=\chiptext] at ({\chipw/2+(\i-1)*\chipw},0) {\i\ \lab};
  }
  \node[circle, fill=courseblue, text=white, font=\Huge\bfseries, minimum size=1.7cm, inner sep=0] at (1.0,-1.7) {#1};
  \node[anchor=west, font=\huge\bfseries, text=courseblue, text width=8.8cm, align=left] at (2.2,-1.7) {#2};
  \node[circle, fill=coursegold!16, draw=coursegold, thick, text=coursegold!75!black, font=\fontsize{26}{26}\selectfont, minimum size=1.9cm, inner sep=0] at (12.6,-1.7) {\faIcon{#3}};
  \ifx\relax#4\relax\else
    \node[anchor=north west, font=\small\itshape, text=coursegray, text width=13.2cm, align=left] at (0.1,-2.95) {Previously: #4};
  \fi
  \node[anchor=north west, font=\footnotesize\bfseries, text=coursegold!65!black] at (0.1,-3.95) {IN THIS PART};
  \node[anchor=north west, font=\small, text width=13.2cm, align=left] at (0.1,-4.35) {#5};
  \end{tikzpicture}
  \end{frame}
}

%
% Before the first \sectiondivider, a deck must define
%   \sectioncount   the number of sections, e.g. 4
%   \sectionlabels  a \foreach list of number/label pairs, e.g. 1/Need,2/Testable,3/Manage,4/Practice
% Icons come from the fontawesome5 package (Font Awesome 5 Free), e.g. \faIcon{robot}.

\usepackage{colortbl}
\usepackage{fontawesome5}
\usetikzlibrary{backgrounds}

% Slides read better without automatic hyphenation.
\hyphenpenalty=10000
\exhyphenpenalty=10000

\definecolor{coursered}{HTML}{A23B2A}
\newcommand{\goldink}{\color{coursegold!65!black}}

\tikzset{
  tile/.style={rounded corners=4pt, fill=courseice, draw=none, align=center, font=\footnotesize, inner sep=5pt},
  tileblue/.style={tile, fill=white, draw=courseblue!55, thick},
  tilegold/.style={tile, fill=coursegold!14, draw=coursegold, thick},
  tilegreen/.style={tile, fill=coursegreen!10, draw=coursegreen!70, thick},
  tilered/.style={tile, fill=coursered!8, draw=coursered!70, thick},
  tilegray/.style={tile, fill=white, draw=coursegray!55, thick, dashed},
  tag/.style={rounded corners=2pt, fill=coursegold, text=white, font=\bfseries\tiny, inner sep=2.5pt},
  chip/.style={rounded corners=3pt, fill=courseblue, text=white, font=\footnotesize\bfseries, align=center, inner sep=4pt},
  note/.style={font=\scriptsize, text=coursegray, align=left},
  arrowgold/.style={-{Latex[length=2.5mm,width=2mm]}, thick, draw=coursegold},
  arrowred/.style={-{Latex[length=2.5mm,width=2mm]}, thick, draw=coursered},
  arrowgreen/.style={-{Latex[length=2.5mm,width=2mm]}, thick, draw=coursegreen},
  arrowgray/.style={-{Latex[length=2mm,width=1.6mm]}, thick, draw=coursegray!70},
}

% A round icon badge, drawn inside a tikzpicture: x y colour icon
\newcommand{\badge}[4]{%
  \node[circle, fill=#3, text=white, minimum size=0.85cm, inner sep=0pt] at (#1,#2) {\large\faIcon{#4}};
}

% Inline tags for the kind of requirement.
\newcommand{\Ftag}{\colorbox{courseblue}{\color{white}\bfseries\scriptsize\ \ F\ \ }}
\newcommand{\NFtag}{\colorbox{coursegold}{\color{white}\bfseries\scriptsize\ NF\ }}
\newcommand{\Dtag}{\colorbox{coursegreen}{\color{white}\bfseries\scriptsize\ \ D\ \ }}
\newcommand{\MLtag}{\colorbox{coursered}{\color{white}\bfseries\scriptsize\ ML\ }}

% Footer: a strip with the concepts that appear for the first time on the slide, above the
% usual bar. The bar total counts only the main lecture; a backup slide is numbered separately.
\makeatletter
\def\conceptstext{}
\def\currentsectionname{}
\setbeamercolor{conceptbar}{bg=coursegold!12,fg=black}
\newcommand{\cterm}[2]{{\bfseries\color{courseblue}#1:}\ #2\hspace{0.8em}}
\newcommand{\newtag}{\colorbox{coursegold}{\color{white}\bfseries\tiny\ NEW\ }\hspace{0.6em}}
\setbeamertemplate{footline}{
  \leavevmode
  \vbox{%
    \ifx\conceptstext\empty\else
      \hbox{\color{coursegold}\rule{\paperwidth}{0.8pt}}%
      \hbox{\begin{beamercolorbox}[wd=\paperwidth,leftskip=0.5cm,rightskip=0.5cm,sep=3pt]{conceptbar}\scriptsize\newtag\conceptstext\end{beamercolorbox}}%
    \fi
    \hbox{%
      \begin{beamercolorbox}[wd=.33\paperwidth,ht=2.8ex,dp=1.2ex,leftskip=0.8em]{palette primary}%
        \ifbeamer@inappendix Backup\else\ifnum\value{section}>0 \arabic{section}\ \textbar\ \currentsectionname\else University of Trieste\fi\fi
      \end{beamercolorbox}%
      \begin{beamercolorbox}[wd=.49\paperwidth,ht=2.8ex,dp=1.2ex,center]{palette tertiary}%
        \coursename{} \textbar{} \coursecode
      \end{beamercolorbox}%
      \begin{beamercolorbox}[wd=.18\paperwidth,ht=2.8ex,dp=1.2ex,right,rightskip=0.8em]{palette secondary}%
        \ifbeamer@inappendix Backup \insertframenumberinappendix\else\insertframenumber{} / \insertmainframenumber\fi
      \end{beamercolorbox}%
    }}}
\makeatother
\newenvironment{cframe}[2]{\gdef\conceptstext{#2}\begin{frame}{#1}}{\end{frame}\gdef\conceptstext{}}

% Section divider: progress strip, recap of the previous part, and the roadmap of the next one.
% Arguments: number, title, icon, previously, items.
\newcommand{\sectiondivider}[5]{%
  \section{#2}\gdef\currentsectionname{#2}
  \begin{frame}
  \centering
  \begin{tikzpicture}
  \pgfmathsetmacro{\chipw}{13.6/\sectioncount}
  \foreach \i/\lab in \sectionlabels {
    \ifnum\i<#1 \def\chipfill{courseblue!22}\def\chiptext{courseblue}\else
    \ifnum\i=#1 \def\chipfill{courseblue}\def\chiptext{white}\else
    \def\chipfill{courseice}\def\chiptext{coursegray}\fi\fi
    \node[rounded corners=3pt, minimum width={(\chipw-0.1)*1cm}, minimum height=0.6cm, font=\scriptsize\bfseries, fill=\chipfill, text=\chiptext] at ({\chipw/2+(\i-1)*\chipw},0) {\i\ \lab};
  }
  \node[circle, fill=courseblue, text=white, font=\Huge\bfseries, minimum size=1.7cm, inner sep=0] at (1.0,-1.7) {#1};
  \node[anchor=west, font=\huge\bfseries, text=courseblue, text width=8.8cm, align=left] at (2.2,-1.7) {#2};
  \node[circle, fill=coursegold!16, draw=coursegold, thick, text=coursegold!75!black, font=\fontsize{26}{26}\selectfont, minimum size=1.9cm, inner sep=0] at (12.6,-1.7) {\faIcon{#3}};
  \ifx\relax#4\relax\else
    \node[anchor=north west, font=\small\itshape, text=coursegray, text width=13.2cm, align=left] at (0.1,-2.95) {Previously: #4};
  \fi
  \node[anchor=north west, font=\footnotesize\bfseries, text=coursegold!65!black] at (0.1,-3.95) {IN THIS PART};
  \node[anchor=north west, font=\small, text width=13.2cm, align=left] at (0.1,-4.35) {#5};
  \end{tikzpicture}
  \end{frame}
}

%
% Before the first \sectiondivider, a deck must define
%   \sectioncount   the number of sections, e.g. 4
%   \sectionlabels  a \foreach list of number/label pairs, e.g. 1/Need,2/Testable,3/Manage,4/Practice
% Icons come from the fontawesome5 package (Font Awesome 5 Free), e.g. \faIcon{robot}.

\usepackage{colortbl}
\usepackage{fontawesome5}
\usetikzlibrary{backgrounds}

% Slides read better without automatic hyphenation.
\hyphenpenalty=10000
\exhyphenpenalty=10000

\definecolor{coursered}{HTML}{A23B2A}
\newcommand{\goldink}{\color{coursegold!65!black}}

\tikzset{
  tile/.style={rounded corners=4pt, fill=courseice, draw=none, align=center, font=\footnotesize, inner sep=5pt},
  tileblue/.style={tile, fill=white, draw=courseblue!55, thick},
  tilegold/.style={tile, fill=coursegold!14, draw=coursegold, thick},
  tilegreen/.style={tile, fill=coursegreen!10, draw=coursegreen!70, thick},
  tilered/.style={tile, fill=coursered!8, draw=coursered!70, thick},
  tilegray/.style={tile, fill=white, draw=coursegray!55, thick, dashed},
  tag/.style={rounded corners=2pt, fill=coursegold, text=white, font=\bfseries\tiny, inner sep=2.5pt},
  chip/.style={rounded corners=3pt, fill=courseblue, text=white, font=\footnotesize\bfseries, align=center, inner sep=4pt},
  note/.style={font=\scriptsize, text=coursegray, align=left},
  arrowgold/.style={-{Latex[length=2.5mm,width=2mm]}, thick, draw=coursegold},
  arrowred/.style={-{Latex[length=2.5mm,width=2mm]}, thick, draw=coursered},
  arrowgreen/.style={-{Latex[length=2.5mm,width=2mm]}, thick, draw=coursegreen},
  arrowgray/.style={-{Latex[length=2mm,width=1.6mm]}, thick, draw=coursegray!70},
}

% A round icon badge, drawn inside a tikzpicture: x y colour icon
\newcommand{\badge}[4]{%
  \node[circle, fill=#3, text=white, minimum size=0.85cm, inner sep=0pt] at (#1,#2) {\large\faIcon{#4}};
}

% Inline tags for the kind of requirement.
\newcommand{\Ftag}{\colorbox{courseblue}{\color{white}\bfseries\scriptsize\ \ F\ \ }}
\newcommand{\NFtag}{\colorbox{coursegold}{\color{white}\bfseries\scriptsize\ NF\ }}
\newcommand{\Dtag}{\colorbox{coursegreen}{\color{white}\bfseries\scriptsize\ \ D\ \ }}
\newcommand{\MLtag}{\colorbox{coursered}{\color{white}\bfseries\scriptsize\ ML\ }}

% Footer: a strip with the concepts that appear for the first time on the slide, above the
% usual bar. The bar total counts only the main lecture; a backup slide is numbered separately.
\makeatletter
\def\conceptstext{}
\def\currentsectionname{}
\setbeamercolor{conceptbar}{bg=coursegold!12,fg=black}
\newcommand{\cterm}[2]{{\bfseries\color{courseblue}#1:}\ #2\hspace{0.8em}}
\newcommand{\newtag}{\colorbox{coursegold}{\color{white}\bfseries\tiny\ NEW\ }\hspace{0.6em}}
\setbeamertemplate{footline}{
  \leavevmode
  \vbox{%
    \ifx\conceptstext\empty\else
      \hbox{\color{coursegold}\rule{\paperwidth}{0.8pt}}%
      \hbox{\begin{beamercolorbox}[wd=\paperwidth,leftskip=0.5cm,rightskip=0.5cm,sep=3pt]{conceptbar}\scriptsize\newtag\conceptstext\end{beamercolorbox}}%
    \fi
    \hbox{%
      \begin{beamercolorbox}[wd=.33\paperwidth,ht=2.8ex,dp=1.2ex,leftskip=0.8em]{palette primary}%
        \ifbeamer@inappendix Backup\else\ifnum\value{section}>0 \arabic{section}\ \textbar\ \currentsectionname\else University of Trieste\fi\fi
      \end{beamercolorbox}%
      \begin{beamercolorbox}[wd=.49\paperwidth,ht=2.8ex,dp=1.2ex,center]{palette tertiary}%
        \coursename{} \textbar{} \coursecode
      \end{beamercolorbox}%
      \begin{beamercolorbox}[wd=.18\paperwidth,ht=2.8ex,dp=1.2ex,right,rightskip=0.8em]{palette secondary}%
        \ifbeamer@inappendix Backup \insertframenumberinappendix\else\insertframenumber{} / \insertmainframenumber\fi
      \end{beamercolorbox}%
    }}}
\makeatother
\newenvironment{cframe}[2]{\gdef\conceptstext{#2}\begin{frame}{#1}}{\end{frame}\gdef\conceptstext{}}

% Section divider: progress strip, recap of the previous part, and the roadmap of the next one.
% Arguments: number, title, icon, previously, items.
\newcommand{\sectiondivider}[5]{%
  \section{#2}\gdef\currentsectionname{#2}
  \begin{frame}
  \centering
  \begin{tikzpicture}
  \pgfmathsetmacro{\chipw}{13.6/\sectioncount}
  \foreach \i/\lab in \sectionlabels {
    \ifnum\i<#1 \def\chipfill{courseblue!22}\def\chiptext{courseblue}\else
    \ifnum\i=#1 \def\chipfill{courseblue}\def\chiptext{white}\else
    \def\chipfill{courseice}\def\chiptext{coursegray}\fi\fi
    \node[rounded corners=3pt, minimum width={(\chipw-0.1)*1cm}, minimum height=0.6cm, font=\scriptsize\bfseries, fill=\chipfill, text=\chiptext] at ({\chipw/2+(\i-1)*\chipw},0) {\i\ \lab};
  }
  \node[circle, fill=courseblue, text=white, font=\Huge\bfseries, minimum size=1.7cm, inner sep=0] at (1.0,-1.7) {#1};
  \node[anchor=west, font=\huge\bfseries, text=courseblue, text width=8.8cm, align=left] at (2.2,-1.7) {#2};
  \node[circle, fill=coursegold!16, draw=coursegold, thick, text=coursegold!75!black, font=\fontsize{26}{26}\selectfont, minimum size=1.9cm, inner sep=0] at (12.6,-1.7) {\faIcon{#3}};
  \ifx\relax#4\relax\else
    \node[anchor=north west, font=\small\itshape, text=coursegray, text width=13.2cm, align=left] at (0.1,-2.95) {Previously: #4};
  \fi
  \node[anchor=north west, font=\footnotesize\bfseries, text=coursegold!65!black] at (0.1,-3.95) {IN THIS PART};
  \node[anchor=north west, font=\small, text width=13.2cm, align=left] at (0.1,-4.35) {#5};
  \end{tikzpicture}
  \end{frame}
}
